Definiamo quindi le seguenti operazioni su $\mathbb{R}^n$

\begin{definition}
    \textbf{Somma in $\mathbb{R}^n$}
    
    \[ \overline{x} + \overline{y} := \left(x_1+y_1, x_2+y_2, \ldots, x_n+y_n\right). \]
    $\mathbb{R}^n$ risulta quindi chiuso rispetto alla somma dato che $\overline{x} +
    \overline{y}\in\mathbb{R}^n$
\end{definition}
\begin{definition}
    \textbf{Prodotto per uno scalare}
    
    \[ \forall \alpha\in\mathbb{R},\; \alpha\cdot\overline{x} :=\left(\alpha x_1,
    \alpha x_2, \ldots, \alpha x_n\right). \]
    $\mathbb{R}^n$ risulta quindi chiuso rispetto al prodotto per uno scalare dato che
    $\forall\alpha\in\mathbb{R}, \; \alpha\cdot\overline{x}\in\mathbb{R}^n$
\end{definition}
\begin{definition}
    \textbf{Prodotto scalare euclideo}
    
    \[
        \langle \overline{x}, \overline{y}\rangle := x_1 y_1 + x_2 y_2 + \ldots + x_n+y_n
         = \sum_{i = 1}^{n} x_i y_i.
    \]
    Indicato anche come
    \[
        \langle \overline{x}, \overline{y}\rangle = \left( \overline{x}, \overline{y} \right)
        = \overline{x}\cdot\overline{y}.
    \]
    Risulta quindi sempre vero che $\langle\overline{x},\overline{y}\rangle\in\mathbb{R}^n$.
\end{definition}
\begin{definition}
    \textbf{Norma euclidea}
    
    \[
        \|\overline{x}\| = \sqrt{\langle \overline{x}, \overline{x}\rangle} =
        \sqrt{x_1 y_1 + x_2 y_2 + \ldots + x_n+y_n} = \sqrt{\sum_{i = 1}^{n} x_i y_i}.
    \]
    Risulta quindi sempre vero che $\|\overline{x}\|\in\mathbb{R}^n$.
\end{definition}